\newtoggle{withbonus} \toggletrue{withbonus} % set true to show bonus points in grade table

% setting underlying course name (Grundlagenfach Mathematik, §-Unterricht, \ldots)
\newcommand{\mycourse}{Ergänzungsfach Informatik}
% name of the class
\newcommand{\myclass}{4. Klassen}
% set date
\newcommand{\mydate}{01. Februar 2024}

\title{\textbf{reguläre und nicht reguläre Sprachen, endliche Automaten}}
\author{Thomas Graf\\
\mycourse, \myclass}
\date{\mydate}
\maketitle
\thispagestyle{headandfoot}		

\begin{center}
	\fbox{\parbox{140mm}{\centering
			\begin{itemize}
				\item Dauer der Prüfung: $40$ Minuten
				%\item \textbf{Zeige in jeder Aufgabe unbedingt Deine Schritte!}
				\item Erlaubte Hilfsmittel:
				\begin{itemize}
				    \item Schreibzeug
				    %\item Taschenrechner
				\end{itemize}
	\end{itemize}}}
\end{center}
\vspace{10mm}

\begin{center}
	Vorname: \rule{45mm}{0.8pt}\qquad Nachname: \rule{45mm}{0.8pt}
\end{center}
\vspace{20mm}

\begin{center}
	\iftoggle{withbonus} {
		\combinedgradetable[h][questions]
	} {
		\gradetable[h][questions]
	}
\end{center}
\vspace{15mm}
\begin{center}
	Note: 
	\vspace{10mm}
	
	\rule{8mm}{1.2pt}
\end{center}
%\thispagestyle{empty}
\clearpage

\begin{questions}

\question{\textbf{Endlichen Automaten zeichnen 1}}

Gegeben ist die Sprache
\begin{align*}
  L_1 := \setcm{x00110y}{x,y\in\lrc{0,1}^*}.  
\end{align*}
\begin{parts}
  \part[1] Nenne ein kürzestes Wort in $L_1$.
    \begin{solution}
      Das Wort $00110$ ist das (eindeutige) kürzeste Wort in $L_1$.
    \end{solution}
  \vspace{5cm}
  \part[3] Entwerfe einen endlichen Automaten (in Diagrammform), welcher $L_1$ akzeptiert.
    \begin{solution}
      \begin{figure}[H]
      \centering
      \begin{tikzpicture}[->,>=stealth',shorten >=1pt,node distance=2.2cm,semithick]
        \tikzstyle{every state}=[fill=blue!10,draw=black,text=black,minimum size=1.4cm]

        \node[initial,state] (qlambda) {$q_{\lambda}$};
        \node[state] (q0) [right of=qlambda] {$q_0$};24
        \node[state] (q00) [right of=q0] {$q_{00}$};
        \node[state] (q001) [right of=q00] {$q_{001}$};
        \node[state] (q0011) [right of=q001] {$q_{0011}$};
        \node[state, accepting] (q00110) [right of=q0011] {$q_{00110}$};

        \path (qlambda) edge [loop above] node {$1$} (qlambda)
              (qlambda) edge [below] node {$0$} (q0)
              
              (q0) edge [below] node {$0$} (q00)
              (q0) edge [above, bend right] node {$1$} (qlambda)
              
              (q00) edge [below] node {$1$} (q001)
              (q00) edge [loop above] node {$0$} (q00)
              
              (q001) edge [below, bend left] node {$0$} (q0)
              (q001) edge [below] node {$1$} (q0011)

              (q0011) edge [below, bend left] node {$1$} (qlambda)
              (q0011) edge [below] node {$0$} (q00110)

              (q00110) edge [loop above] node {$0, 1$} (q00110);
      \end{tikzpicture}
      \caption{Endlicher Automat (in Diagrammform) für die Sprache $L_1$}
      \label{fig:L1}
      \end{figure}
    \end{solution}
\end{parts}
\droptotalpoints
\clearpage \question{\textbf{Endlichen Automaten zeichnen 2}} Gegeben ist die Sprache \begin{align*} L_2 := \setcm{x1y}{x,y\in\lrc{0,1}^*, \abs{y}\geq 1}.  
\end{align*}
\begin{parts}
  \part[1] Nenne ein kürzestes Wort in $L_2$.
    \begin{solution}
      Die beiden Wörter $10$ und $11$ sind die beiden kürzesten Wörter in $L_2$.
    \end{solution}
  \vspace{5cm}
  \part[3] Entwerfe einen endlichen Automaten (in Diagrammform), welcher $L_2$ akzeptiert.
    \begin{solution}
      \begin{figure}[H]
      \centering
      \begin{tikzpicture}[->,>=stealth',shorten >=1pt,node distance=3.2cm,semithick]
        \tikzstyle{every state}=[fill=blue!10,draw=black,text=black,minimum size=1.4cm]

        \node[initial,state] (qlambda) {$q_{\lambda}$};
        \node[state] (q1) [right of=qlambda] {$q_1$};
        \node[state,accepting] (q2) [right of=q1] {$q_{2}$};

        \path (qlambda) edge [loop above] node {$0$} (qlambda)
              (qlambda) edge [below] node {$1$} (q1)
              
              (q1) edge [below] node {$0, 1$} (q2)

              (q2) edge [loop above] node {$0, 1$} (q2);
      \end{tikzpicture}
      \caption{Endlicher Automat (in Diagrammform) für die Sprache $L_2$}
      \label{fig:L2}
      \end{figure}
    \end{solution}
\end{parts}
\droptotalpoints
\clearpage


\question{\textbf{\green{Reguläre} Sprache}}

Gegeben ist die Sprache
\begin{align*}
  L_3 := \setcm{0^i1^j}{i,j\in\N}.  
\end{align*}
\begin{parts}
  \part[1] Nenne genau drei verschiedene Wörter, die in $L_3$ liegen.
    \begin{solution}
      $\lambda, 0, 1, 00, 00111, 011, 000001$ sind sieben Beispiele von verschiedenen Wörtern aus $L_3$. 
    \end{solution}
  \vspace{5cm}
  \part[3] Beweise, dass $L_3$ \green{regulär} ist.
    \begin{solution}
      Um zu beweisen, dass $L_3$ regulär ist, genügt es einen endlichen Automaten für $L_3$ anzugeben. Dieser endliche
      Automat ist in \cref{fig:L3} angegeben.
      \begin{figure}[H]
      \centering
      \begin{tikzpicture}[->,>=stealth',shorten >=1pt,node distance=3.2cm,semithick]
        \tikzstyle{every state}=[fill=blue!10,draw=black,text=black,minimum size=1.4cm]

        \node[initial,state,accepting] (q0) {$q_{0}$};
        \node[state,accepting] (q01) [right of=q0] {$q_{01}$};
        \node[state] (qt) [right of=q01] {$q_{t}$};

        \path (q0) edge [loop above] node {$0$} (q0)
              (q0) edge [below] node {$1$} (q01)
              
              (q01) edge [loop above] node {$1$} (q01)
              (q01) edge [below] node {$1$} (qt)

              (qt) edge [loop above] node {$0, 1$} (qt);
      \end{tikzpicture}
      \caption{Endlicher Automat (in Diagrammform) für die Sprache $L_3$}
      \label{fig:L3}
      \end{figure}
    \end{solution}
\end{parts}
\droptotalpoints
\clearpage


\question{\textbf{\red{Nicht reguläre} Sprache}}

Gegeben ist die Sprache
\begin{align*}
  L_4 := \setcm{w\#w}{w\in\lrc{a,b}^*}.  
\end{align*}
\begin{parts}
  \part[1] Nenne genau drei verschiedene Wörter, die in $L_4$ liegen. 
  \begin{solution}
    \begin{align*}
     \#, a\#a, bbabaaaab\#bbabaaaab
    \end{align*} 
  \end{solution}
  \vspace{5cm}
  \part[3] Zeige, dass $L_4$ \red{nicht regulär} ist.
  \begin{solution}
    Angenommen, $L_4$ sei regulär. Dann existiert ein endlicher Automat
    \begin{align*}
      A = (Q, {a,b}, \delta, q_0, F) 
    \end{align*}
    und $A$ akzeptiert $L_4$. Sei nun $m := \abs{Q}$, also die Anzahl der verschiedenen Zustände des endlichen Automaten
    $A$. Wir betrachten die $m+1$ verschiedenen Wörter
    \begin{align*}
      a^1\#, a^2\#, \ldots, a^{m+1}\#.
    \end{align*}
    Dies sind mehr Worte (genau ein Wort mehr), als $A$ Zustände hat. Nach dem Schubfachprinzip (pigeonhole principle) existieren
    $i,j\in\lrc{1,2,\ldots,m+1}$ mit $i<j$, sodass
    \begin{align*}
      \hat{\delta}\lr{q_0,a^i\#} = \hat{\delta}\lr{q_0,a^j\#}. 
    \end{align*}
    Nach der im Unterricht bewiesenen Aussage gilt
    \begin{align*}
      a^i\#\in L_4 \iff a^j\#\in L_4 
    \end{align*}
    für alle $z\in\lrc{a,b}^*$. Doch dies gilt nicht, da zum Beispiel für die Wahl $z:=a^i$
    \begin{align*}
      a^i\# z = a^i\# a^i\in L_4 \quad\text{aber}\quad a^j\# z = a^j\# a^i\notin L_4
    \end{align*}
    gilt, da $i<j$ und somit insbesondere $i\neq j$. Damit war unsere ursprüngliche Annahme falsch und es existiert kein
    endlicher Automat für die Sprache $L_4$.
  \end{solution}
\end{parts}
\droptotalpoints

\end{questions}

